% Bibliographical Notes for Chapter II
% Chandrasekhar, Hydrodynamic and Hydromagnetic Stability
% Chapter II: The Thermal Instability of a Layer of Fluid Heated from Below: 1. The B\'enard Problem

\section*{BIBLIOGRAPHICAL NOTES}

The phenomenon of thermal convection in fluids, as we now understand it, was
discovered by Count Rumford:

\begin{enumerate}

\item \textsc{Count Rumford}, ``Of the propagation of heat in fluids'',
  \textit{Complete Works}, \textbf{1}, 239, American Academy of Arts and Sciences,
  Boston, 1870.

\end{enumerate}

\noindent An interesting account of Rumford's discovery with relevant extracts from his
writings will be found in:

\begin{enumerate}
\setcounter{enumi}{1}

\item S.~C. \textsc{Brown}, ``Count Rumford discovers thermal convection'',
  \textit{Daedalus}, \textbf{86}, 340--3 (1957).

\end{enumerate}

\noindent In reference~2, Brown points out that the term `convection' was first introduced
in science by William Prout:

\begin{enumerate}
\setcounter{enumi}{2}

\item W. \textsc{Prout}, \textit{Bridgewater Treatises}, \textbf{8}, 65, edited
  by W.~Pickering, London, 1834.

\end{enumerate}

\noindent The following extract from Prout is taken from reference~2:

\begin{quote}
`There is at present no single term in our language employed to denote this
mode of propagation of heat; but we venture to propose for that purpose the
term \textit{convection} (\textit{convectio}, a carrying or converging), which
not only expresses the leading facts, but also accords very well with the two
other terms (conduction and radiation).'
\end{quote}

\noindent The fact that convection in horizontal layers of fluid is often accompanied by
a cellular pattern of motions seems to have been first observed by James Thomson:

\begin{enumerate}
\setcounter{enumi}{3}

\item J. \textsc{Thomson}, ``On a changing tesselated structure in certain liquids'',
  \textit{Proc.\ Phil.\ Soc.\ Glasgow}, \textbf{13}, 464--8 (1882).

\end{enumerate}

\noindent The following extract is from reference~4:

\begin{quote}
`The occurrence of the phenomena seems to be associated essentially with
cooling of the liquid at its surface where exposed to the air, when the main
body of the liquid is at a temperature somewhat above that assumed by a thin
superficial film. A very slight excess of temperature in the body of the fluid
above that of the surrounding air is sufficient to institute the tesselated
changing structure\,\ldots. The motions now described and explained as occurring
in the soapy water and other liquids when showing the tesselated structure
he [Thomson] thus believes constitute one special case of what is often called
convective circulation. He points out, however, that this case, with a thin
film cooled at the surface while the great body of the liquid is at a somewhat
higher temperature, presents essential distinctions from the case in which
convective circulation is caused by heat applied at the bottom of the vessel.
The surface phenomena and the actual motions throughout the body of the
liquid are very different in the two cases, as may easily be seen by a little
consideration, both theoretical and observational.'
\end{quote}

\noindent The first quantitative experiments on the onset of thermal instability and the
recognition of the role of viscosity in the phenomenon are due to H.~B\'enard:

\begin{enumerate}
\setcounter{enumi}{4}

\item H. \textsc{B\'enard}, ``Les Tourbillons cellulaires dans une nappe liquide'',
  \textit{Revue g\'en\'erale des Sciences pures et appliqu\'ees}, \textbf{11},
  1261--71 and 1309--28 (1900).

\item -------, ``Les Tourbillons cellulaires dans une nappe liquide transportant de
  la chaleur par convection en r\'egime permanent'', \textit{Annales de Chimie et de
  Physique}, \textbf{23}, 62--144 (1901).

\end{enumerate}

\noindent On the theoretical side, the fundamental paper is that of Lord Rayleigh:

\begin{enumerate}
\setcounter{enumi}{6}

\item \textsc{Lord Rayleigh}, ``On convective currents in a horizontal layer of fluid
  when the higher temperature is on the under side'', \textit{Phil.\ Mag.},
  \textbf{32}, 529--46 (1916); also \textit{Scientific Papers}, \textbf{6}, 432--46,
  Cambridge, England, 1920.

\end{enumerate}

\noindent The paper by Lord Rayleigh is included in an anthology of related papers edited
by Saltzmann:

\begin{enumerate}
\setcounter{enumi}{7}

\item B. \textsc{Saltzmann}, ``The general circulation as a problem in thermal convection;
  A collection of classical and modern theoretical papers'', \textit{Scientific Report
  No.~1}, General Circulation Project, Department of Meteorology, Massachusetts
  Institute of Technology, 1958.

\end{enumerate}

\noindent A collection of papers containing some beautiful illustrations of convection
phenomena in Nature is:

\begin{enumerate}
\setcounter{enumi}{8}

\item \textit{Convection Patterns in the Atmosphere and Ocean}, edited by
  R.~W.~Miner, \textit{Annals of the New York Academy of Sciences},
  \textbf{48}, 1947.

\end{enumerate}

\noindent References to the B\'enard convection cells from the standpoint of meteorology
will be found in:

\begin{enumerate}
\setcounter{enumi}{9}

\item D. \textsc{Brunt}, \textit{Physical and Dynamical Meteorology}, 219--21,
  Cambridge, England, 1939.

\item O.~G. \textsc{Sutton}, \textit{Micrometeorology}, 119--25, McGraw-Hill Book Co.\
  Inc., New York, 1953.

\end{enumerate}

\noindent And a general account of the subject treated in this and the following chapters
is given in:

\begin{enumerate}
\setcounter{enumi}{11}

\item S. \textsc{Chandrasekhar}, ``Thermal convection'', \textit{Daedalus},
  \textbf{86}, 323--39 (1957).

\end{enumerate}

\noindent \S\,7. A useful reference for the foundations of hydrodynamics is:

\begin{enumerate}
\setcounter{enumi}{12}

\item S. \textsc{Goldstein}, \textit{Modern Developments in Fluid Dynamics},
  \textbf{1}, 90--101 and \textbf{2}, 601--9, Oxford, England, 1938.

\end{enumerate}

\noindent See also:

\begin{enumerate}
\setcounter{enumi}{13}

\item W.~F. \textsc{Cope}, ``The equations of hydrodynamics in a very general form'',
  \textit{Aeronautical Research Committee Reports and Memoranda}, No.~1903, 1--6,
  (1942).

\end{enumerate}

\noindent For the notation of Cartesian tensors see:

\begin{enumerate}
\setcounter{enumi}{14}

\item H. \textsc{Jeffreys}, \textit{Cartesian Tensors}, Cambridge, England, 1931.

\item H. and B.~S. \textsc{Jeffreys}, \textit{Methods of Mathematical Physics},
  chapter~3, Cambridge, England, 1956.

\end{enumerate}

\noindent \S\,8. The approximation introduced in this section is due to:

\begin{enumerate}
\setcounter{enumi}{16}

\item J. \textsc{Boussinesq}, \textit{Th\'eorie Analytique de la Chaleur}, \textbf{2},
  172, Gauthier-Villars, Paris, 1903.

\end{enumerate}

\noindent \S\S\,9 and~10. As we have stated, the foundations of the subject were laid by Lord
Rayleigh in reference~7. The theory, from a somewhat more general point of
view and leading to equation~(104), is due to Jeffreys:

\begin{enumerate}
\setcounter{enumi}{17}

\item H. \textsc{Jeffreys}, ``The stability of a layer of fluid heated below'',
  \textit{Phil.\ Mag.}, \textbf{2}, 833--44 (1926).

\item -------, ``Some cases of instability in fluid motion'',
  \textit{Proc.\ Roy.\ Soc.\ (London)} A, \textbf{118}, 195--208 (1928).

\end{enumerate}

\noindent See also:

\begin{enumerate}
\setcounter{enumi}{19}

\item H. \textsc{Jeffreys}, ``The instability of a compressible fluid heated below'',
  \textit{Proc.\ Camb.\ Phil.\ Soc.}, \textbf{26}, 170--2 (1930).

\end{enumerate}

\noindent In reference~20, Jeffreys shows that the criterion for the onset of instability
derived for incompressible fluids can, under certain conditions, be used for
compressible fluids if $\beta$ is interpreted as the difference between the adiabatic
temperature gradient and the prevailing temperature gradient.

\noindent \S\,11. The validity of the principle of the exchange of stabilities for this problem
was proved by Rayleigh (reference~7) for the case of two free boundaries. The
proof in the general case is due to Pellew and Southwell:

\begin{enumerate}
\setcounter{enumi}{20}

\item A. \textsc{Pellew} and R.~V. \textsc{Southwell}, ``On maintained convective motion in
  a fluid heated from below'', \textit{Proc.\ Roy.\ Soc.\ (London)} A, \textbf{176},
  312--43 (1940).

\end{enumerate}

\noindent \S\,13. The first of the two variational principles considered in this section is due
to Pellew and Southwell (reference~21); for the second of these see:

\begin{enumerate}
\setcounter{enumi}{21}

\item S. \textsc{Chandrasekhar}, ``On characteristic value problems in high order
  differential equations which arise in studies on hydrodynamic and hydromagnetic
  stability'', \textit{American Math.\ Monthly}, \textbf{61}, 32--45 (1954).

\end{enumerate}

\noindent The fact that the variational principle gives the true minimum does not seem to
have been proved before.

\noindent \S\,14. The thermodynamic significance of the variational principle was first
considered explicitly though in the context of a somewhat different problem by:

\begin{enumerate}
\setcounter{enumi}{22}

\item H. \textsc{Jeffreys}, ``The thermodynamics of thermal instability in liquids'',
  \textit{Quart.\ J.\ Mech.\ Appl.\ Math.}, \textbf{9}, 1--5 (1956).

\end{enumerate}

\noindent Less explicitly, it is contained in:

\begin{enumerate}
\setcounter{enumi}{23}

\item W.~V.~R. \textsc{Malkus}, ``The heat transport and spectrum of thermal turbulence'',
  \textit{Proc.\ Roy.\ Soc.\ (London)} A, \textbf{225}, 196--212 (1954).

\end{enumerate}

\noindent A complete discussion of the underlying principle including sources of dissipation,
besides viscosity, and allowing for overstability will be found in:

\begin{enumerate}
\setcounter{enumi}{24}

\item S. \textsc{Chandrasekhar}, ``The thermodynamics of thermal instability in liquids'',
  \textit{Max Planck Festschrift 1958}, 103--14, Veb Deutscher Verlag der Wissenschaften,
  Berlin, 1958.

\end{enumerate}

\noindent \S\,15. The exact solution of the basic characteristic value problem was first
accomplished by:

\begin{enumerate}
\setcounter{enumi}{25}

\item A.~R. \textsc{Low}, ``On the criterion for stability of a layer of viscous fluid heated
  from below'', \textit{Proc.\ Roy.\ Soc.\ (London)} A, \textbf{125}, 180--95 (1929).

\end{enumerate}

\noindent By a somewhat different method it was also carried out by Pellew and Southwell
(reference~21). The method followed in the text is essentially that of Pellew
and Southwell. The definitive treatment of the problem is due to:

\begin{enumerate}
\setcounter{enumi}{26}

\item W.~H. \textsc{Reid} and D.~L. \textsc{Harris}, ``Some further results on the B\'enard
  problem'', \textit{The Physics of Fluids}, \textbf{1}, 102--10 (1958).

\item -------, ``Streamlines in B\'enard convection cells'', ibid.\ \textbf{2}, 716--17
  (1959).

\end{enumerate}

\noindent \S\,16. The solution of the equation
%
\[ \frac{\partial^2\phi}{\partial x^2} + \frac{\partial^2\phi}{\partial y^2} = -a^2\phi \]
%
appropriate for a hexagonal cell pattern was discovered by:

\begin{enumerate}
\setcounter{enumi}{28}

\item D.~G. \textsc{Christopherson}, ``Note on the vibration of membranes'',
  \textit{Quart.\ J.\ of Math.\ (Oxford Series)}, \textbf{11}, 63--65 (1940).

\end{enumerate}

\noindent The generalization of Christopherson's solution given in \S\,(e) is due to:

\begin{enumerate}
\setcounter{enumi}{29}

\item F.~E. \textsc{Bisshopp}, ``On two-dimensional cell patterns'',
  \textit{J.\ Math.\ Analysis and Applications}, \textbf{1}, 373--85 (1960).

\end{enumerate}

\noindent \S\,17. Variational solutions have been given by Pellew and Southwell (reference~21)
and Reid and Harris~(27).

\noindent \S\,18. The references for the experimental work described in this section are:

\begin{enumerate}
\setcounter{enumi}{30}

\item H. \textsc{B\'enard} and D. \textsc{Avsec}, ``Travaux r\'ecents sur les tourbillons cellulaires
  et les tourbillons en bandes; applications \`a l'astrophysique et \`a la m\'et\'eorologie'',
  \textit{Le Journal de physique et le radium}, \textbf{9}, 486--500 (1938).

\item R.~J. \textsc{Schmidt} and S.~W. \textsc{Milverton}, ``On the instability of a fluid when
  heated from below'', \textit{Proc.\ Roy.\ Soc.\ (London)} A, \textbf{152}, 586--94
  (1935).

\item ------- and O.~A. \textsc{Saunders}, ``On the motion of a fluid heated from below'',
  ibid.\ \textbf{165}, 216--28 (1938).

\item O.~A. \textsc{Saunders}, M. \textsc{Fishenden}, and H.~D. \textsc{Mansion}, ``Some measurements
  of convection by an optical method'', \textit{Engineering}, \textbf{139}, 483--5 (1935).

\item P.~L. \textsc{Silveston}, ``W\"armedurchgang in waagerechten Fl\"ussigkeitsschichten'',
  Part~1, \textit{Forsch.\ Ing.\ Wes.}, \textbf{24}, 29--32 and 59--69 (1958).

\end{enumerate}

\noindent See also:

\begin{enumerate}
\setcounter{enumi}{35}

\item W.~V.~R. \textsc{Malkus}, ``Discrete transitions in turbulent convection'',
  \textit{Proc.\ Roy.\ Soc.\ (London)} A, \textbf{225}, 185--95 (1954).

\end{enumerate}

\noindent In reference~31, B\'enard's earlier work (references~5 and~6) is summarized and
applications to astrophysical and meteorological phenomena are considered. The
role of surface tension in B\'enard's experiments is discussed by:

\begin{enumerate}
\setcounter{enumi}{36}

\item J.~R.~A. \textsc{Pearson}, ``On convection cells induced by surface tension'',
  \textit{J.\ Fluid Mech.}, \textbf{4}, 489--500 (1958).

\end{enumerate}

\noindent The `Schmidt--Milverton' principle is described in reference~32. And the experimental
arrangement used in securing the photographs illustrated in Fig.~15 is
described in reference~34.

\medskip

\noindent Matters closely related to the subject-matter of this chapter, but not explicitly
considered, are treated in:

\begin{enumerate}
\setcounter{enumi}{37}

\item H. \textsc{Tippelskirch}, ``Weitere Konvektionsversuche: der Nachweis der
  Ringzellen und ihrer Verallgemeinerung'', \textit{Beitr\"age zur Physik der
  Atmosph\"are}, \textbf{32}, 2--22 (1959).

\item -------, ``\"Uber Konvektionszellen, insbesondere im fl\"ussigen Schwefel'',
  ibid.\ \textbf{29}, 37--54 (1956).

\item J. \textsc{Zierep}, ``\"Uber die Bevorzugung der Sechseckzellen bei Konvektionsstr\"omungen
  \"uber einer gleichm\"assig erw\"armten Grundfl\"ache'', ibid.\ \textbf{31}, 31--39
  (1958).

\item -------, ``Zur Theorie der Zellularkonvektion III'', ibid.\ \textbf{32}, 23--33 (1959).

\item -------, ``\"Uber rotationssymmetrische Zellularkonvektionsstr\"omungen'',
  \textit{Z.\ f\"ur angewandte Mathematik und Mechanik}, \textbf{38}, 1--4 (1958).

\end{enumerate}

\noindent Some progress has recently been achieved in developing theories which will be
applicable to Rayleigh numbers slightly in excess of $R_c$ and when the amplitudes
are finite. The following references may be noted in this connexion:

\begin{enumerate}
\setcounter{enumi}{42}

\item W.~V.~R. \textsc{Malkus} and G. \textsc{Veronis}, ``Finite amplitude cellular convection'',
  \textit{J.\ Fluid Mech.}, \textbf{4}, 225--60 (1958).

\item L.~P. \textsc{Gor'kov}, ``Stationary convection in a plane liquid layer near the
  critical heat transfer point'', \textit{Soviet Physics, JETP}, \textbf{6}, 311--15 (1958).

\end{enumerate}

\noindent See also Appendix~I.
